%%%%%%%%%%%%%%%%%%%%%%%%%%%%%%%%%%%%%%%%%%%%%%%%%%%%%%%%%%%%%%%%%%
%% Extensions of Schur Numbers
%%%%%%%%%%%%%%%%%%%%%%%%%%%%%%%%%%%%%%%%%%%%%%%%%%%%%%%%%%%%%%%%%%
\documentclass[11pt]{article}

\usepackage[margin=1in]{geometry}
\usepackage{amsmath,amssymb,amsthm}
\usepackage{mathtools}
\usepackage{hyperref}
\usepackage{booktabs}
\usepackage{microtype}
\usepackage{enumitem}

\theoremstyle{plain}
\newtheorem{theorem}{Theorem}[section]
\newtheorem{lemma}[theorem]{Lemma}
\newtheorem{proposition}[theorem]{Proposition}
\newtheorem{corollary}[theorem]{Corollary}
\newtheorem{conjecture}[theorem]{Conjecture}

\theoremstyle{definition}
\newtheorem{definition}[theorem]{Definition}
\newtheorem{remark}[theorem]{Remark}
\newtheorem{observation}[theorem]{Observation}

\newcommand{\ZZ}{\mathbb{Z}}
\newcommand{\DS}{\mathrm{DS}}
\newcommand{\MS}{\mathrm{MS}}

\begin{document}

\title{Extensions of Schur Numbers:\\
Order-Invariance, Density Thresholds, and Multiplicative Variants}

\author{Alex Towell\thanks{Department of Computer Science,
Southern Illinois University Edwardsville.
Email: \texttt{atowell@siue.edu}}}

\date{}
\maketitle

\begin{abstract}
We report computational results in three directions extending
classical Schur numbers.
First, we study Schur numbers $S(G,k)$ for finite abelian groups~$G$
and find that $S(G,1)$ and $S(G,2)$ depend only on $|G|$, verified
through order~20, while $S(G,3)$ is \emph{not} order-invariant:
$S(\ZZ/9\ZZ,3)=8$ but $S(\ZZ/3\ZZ\times\ZZ/3\ZZ,3)=7$.
Second, we introduce density Schur numbers $\DS(k,\alpha)$ and
determine $\DS(2,\tfrac{1}{2})=5$ with a Lean~4 proof,
identifying a phase transition in the parameter~$\alpha$.
Third, we tabulate Rado numbers for several linear equations and
compute multiplicative Schur numbers $\MS(k)$ on the domain
$\{2,\ldots,N\}$, finding $\MS(1)=3$ and $\MS(2)=31$.
All results are computational (exhaustive enumeration or SAT-based)
and should be regarded as such until independent formal verification.
\end{abstract}

%% ================================================================
\section{Schur numbers in finite abelian groups}\label{sec:groups}
%% ================================================================

\subsection{Definitions}

Let $G$ be a finite abelian group written additively.
A subset $C\subseteq G$ is \emph{sum-free} if there are no
$a,b,c\in C$ with $a+b=c$ (the case $a=b$ is permitted, so
$\{0\}$ is never sum-free since $0+0=0$).

\begin{definition}
$S(G,k)$ is the maximum size of a subset $A\subseteq G$ admitting a
partition $A=C_1\sqcup\cdots\sqcup C_k$ into $k$ sum-free sets.
\end{definition}

\noindent
Equivalently, $S(G,k)+1$ is the minimum size of a subset of~$G$ that
forces a monochromatic Schur triple in every $k$-coloring.
When $G=\ZZ/n\ZZ$ and $k=1$, the quantity $S(G,1)$ is the maximum
sum-free subset of~$\ZZ/n\ZZ$.

For a single color ($k=1$), Green and
Ruzsa~\cite{green-ruzsa} established asymptotically tight bounds on
the maximum sum-free subset of a general abelian group.
For $\ZZ/p\ZZ$ with $p$ an odd prime, the interval
$\{x : p/3 < x < 2p/3\}$ is a maximum sum-free set of size
$\lfloor p/3 \rfloor + O(1)$; we have verified this computationally
for all primes $p \le 19$.

\subsection{Order-invariance}

We say that $S(-,k)$ is \emph{order-invariant} if $S(G,k)$
depends only on $|G|$ for all finite abelian groups~$G$.
Since any group isomorphism preserves sum-freeness, $S(G,k)$
is certainly an isomorphism invariant.
The question is whether non-isomorphic groups of the same order can
differ.

\begin{observation}[Computational, verified through order 20]\label{obs:inv}
$S(G,1)$ and $S(G,2)$ are order-invariant for finite abelian groups
of order $\le 20$.
This covers 30 non-isomorphic groups, including all~5 abelian
groups of order~16 (which all satisfy $S(G,1)=8$ and $S(G,2)=12$).
\end{observation}

\noindent
For $k=1$, order-invariance may follow from the Green--Ruzsa
theory~\cite{green-ruzsa}, which determines $S(G,1)$ in terms of
$|G|$ and the group's structure constants.
The $k=2$ invariance appears to be a new observation.

\begin{observation}[Computational]\label{obs:break}
$S(G,3)$ is \emph{not} order-invariant.
The first failure occurs at order~$9$:
\[
S(\ZZ/9\ZZ,\, 3) = 8, \qquad
S(\ZZ/3\ZZ\times\ZZ/3\ZZ,\, 3) = 7.
\]
A second failure occurs at order~12:
$S(\ZZ/3\ZZ\times\ZZ/4\ZZ,\, 3) = 11$ while
$S(\ZZ/2\ZZ\times\ZZ/2\ZZ\times\ZZ/3\ZZ,\, 3) = 10$.
\end{observation}

\noindent
The underlying mechanism appears related to the group exponent:
$\ZZ/9\ZZ$ has exponent~$9$ while
$(\ZZ/3\ZZ)^2$ has exponent~$3$, and the higher exponent allows more
room for three independent sum-free color classes.

\subsection{Tables}

Table~\ref{tab:cyclic} records $S(\ZZ/n\ZZ, k)$ for small~$n$ and
$k=1,2,3$.

\begin{table}[ht]
\centering
\caption{Schur numbers $S(\ZZ/n\ZZ,k)$ for $n=2,\ldots,16$.
All values computed by exhaustive enumeration.}\label{tab:cyclic}
\smallskip
\begin{tabular}{@{}r rrr r@{}}
\toprule
$n$ & $S(-,1)$ & $S(-,2)$ & $S(-,3)$ & $S_2/n$ \\
\midrule
 2  &  1 &  1 &  1 & 0.500 \\
 3  &  1 &  2 &  2 & 0.667 \\
 4  &  2 &  3 &  3 & 0.750 \\
 5  &  2 &  4 &  4 & 0.800 \\
 6  &  3 &  4 &  5 & 0.667 \\
 7  &  2 &  4 &  6 & 0.571 \\
 8  &  4 &  6 &  7 & 0.750 \\
 9  &  3 &  6 &  8 & 0.667 \\
10  &  5 &  8 &  9 & 0.800 \\
11  &  4 &  8 & 10 & 0.727 \\
12  &  6 &  9 &    & 0.750 \\
13  &  4 &  8 &    & 0.615 \\
14  &  7 &  9 &    & 0.643 \\
15  &  6 & 12 &    & 0.800 \\
16  &  8 & 12 &    & 0.750 \\
\bottomrule
\end{tabular}
\end{table}

\noindent
Blank entries in the $k=3$ column indicate that the exhaustive
computation ($3^n$ colorings) was not run for those~$n$ due to time
constraints.
For $k=3$, note that $S(\ZZ/n\ZZ,3) = n-1$ for all $n \le 11$:
the $n-1$ nonzero elements of $\ZZ/n\ZZ$ can always be 3-colored
with each color class sum-free, at least through $n=11$.

\begin{conjecture}
$S(G,2)$ is order-invariant for all finite abelian groups.
\end{conjecture}

%% ================================================================
\section{Density Schur numbers}\label{sec:density}
%% ================================================================

\begin{definition}\label{def:ds}
For integers $k\ge 1$ and a real parameter $\alpha\in[0,1]$,
define the \emph{density Schur number} $\DS(k,\alpha)$ as the
minimum~$N$ such that every $k$-coloring of $\{1,\ldots,N\}$ has at
least one color class~$C_i$ that is \emph{good}: either $C_i$
contains a Schur triple ($a+b=c$ with $a,b,c\in C_i$ and
$a,b,c\ge 1$) or $|C_i| > \alpha N$.
\end{definition}

\begin{remark}[Indexing caveat]\label{rem:index}
Definition~\ref{def:ds} uses the 1-indexed domain $\{1,\ldots,N\}$,
which is the natural setting for integer Schur triples.
An alternative 0-indexed definition on $\{0,\ldots,N-1\}$ would give
different values because $0+0=0$ is a ``free'' Schur triple that
forces any color class containing~0 to be good by default.
Our Lean~4 formalization~\cite{code-repo} uses the
1-indexed convention (elements $\{1,\ldots,N\}$, represented
internally as \texttt{Fin~N} with the Schur condition
$(a+1)+(b+1)=(c+1)$).
\end{remark}

\begin{theorem}[Lean-verified]\label{thm:ds2}
$\DS(2,\tfrac{1}{2}) = 5$.
That is, for every 2-coloring of $\{1,2,3,4,5\}$, some
color class either contains a Schur triple or has more than
$\tfrac{1}{2}\cdot 5 = 2.5$ elements.
The witness coloring on $\{1,2,3,4\}$ that avoids this is
$C_0=\{1,4\}$, $C_1=\{2,3\}$: both classes are sum-free and both
have exactly $2 \le \lfloor 4/2 \rfloor$ elements, so neither
exceeds the density threshold.
\end{theorem}

\noindent
The proof is formalized in Lean~4 with zero sorry statements.
The lower bound (exhibiting the witness coloring) and upper bound
(pigeonhole: some class has $\ge 3 > 5/2$ elements) are both
machine-checked.

\begin{theorem}[Pigeonhole]\label{thm:ds-pigeonhole}
$\DS(k,\, 1/k) = S(k)+1$,
where $S(k)$ is the classical Schur number.
\end{theorem}

\begin{proof}
At $N = S(k)+1$, every $k$-coloring of $\{1,\ldots,N\}$ forces a
monochromatic Schur triple by the definition of~$S(k)$, so every
color class is good (the Schur-triple condition alone suffices).
At $N = S(k)$, there exists a $k$-coloring with no monochromatic
Schur triple.
By pigeonhole, some class has $\ge \lceil N/k \rceil$ elements, but
the density condition requires \emph{strictly} more than $N/k$
elements; when $k$ divides~$N$ and the partition is balanced, this
threshold is not exceeded.
\end{proof}

\subsection{Phase transition for $k=2$}

A related quantity is the \emph{density Ramsey number}
$\DS'(k,\alpha)$: the minimum~$N$ such that every $k$-coloring of
$\{1,\ldots,N\}$ has a color class~$C$ satisfying \emph{both}
$|C|\ge\alpha N$ \emph{and} $C$ contains a Schur triple (the
conjunction, not the disjunction in Definition~\ref{def:ds}).
Computational experiments reveal three regimes for
$\DS'(2,\alpha)$:

\begin{enumerate}[nosep]
\item \textbf{Regime~1} ($\alpha \le 0.59$): $\DS'(2,\alpha) = 5$,
matching $S(2)+1$.
\item \textbf{Regime~2} ($0.60 \le \alpha \le 0.66$):
$\DS'(2,\alpha) = 6$.
\item \textbf{Regime~3} ($\alpha \ge 0.67$): $\DS'(2,\alpha) > 12$.
\end{enumerate}

\noindent
The transition from Regime~1 to Regime~2 is sharp: at $\alpha=0.60$,
$N=5$ has an avoiding coloring (one class can have 3 out of 5
elements, and $3 < 0.60 \times 5 = 3.0$, so the density condition
is not triggered), but $N=6$ does not.
These values are determined by exhaustive enumeration over all
$2^N$ colorings.

%% ================================================================
\section{Rado numbers and multiplicative Schur}\label{sec:rado}
%% ================================================================

\subsection{Rado numbers for linear equations}

A Rado number $R(E,k)$ for an equation~$E$ is the minimum~$N$ such
that every $k$-coloring of $\{1,\ldots,N\}$ has a monochromatic
solution to~$E$.
We computed the following values using a SAT-based approach
(CaDiCaL 1.5 solver with direct encoding).

\begin{table}[ht]
\centering
\caption{Rado numbers $R(E,k)$ for small equations and
$k=1,2,3$.}\label{tab:rado}
\smallskip
\begin{tabular}{@{}l rrr@{}}
\toprule
Equation $E$ & $k=1$ & $k=2$ & $k=3$ \\
\midrule
$a+b=c$ (Schur)       &  2 &  5 & 14 \\
$a+b=2c$ (van der Waerden) &  --- &  9 & 27 \\
$a+b+c=d$             &  --- & 11 & 43 \\
$a+2b=3c$             &  --- & 13 & --- \\
\bottomrule
\end{tabular}
\end{table}

\noindent
The Schur values $R(a{+}b{=}c,\, k) = S(k)+1$ are well-known
($S(1)=1$, $S(2)=4$, $S(3)=13$, $S(4)=44$).
The van der Waerden values match $W(k;3)$.
The values $R(a{+}b{+}c{=}d,\, 2)=11$ and
$R(a{+}b{+}c{=}d,\, 3)=43$ are verified by our SAT solver and are
consistent with the known literature on Rado's theorem.
A dash (---) indicates that the value was not computed (either
trivial or exceeding our search bound).

\subsection{Multiplicative Schur numbers}

\begin{definition}
$\MS(k)$ is the minimum~$N$ such that every $k$-coloring of
$\{2,3,\ldots,N\}$ contains a monochromatic \emph{multiplicative
triple} $\{a,b,ab\}$ (with $a\le b$).
\end{definition}

\noindent
The domain starts at~2 to exclude~1, which would trivially satisfy
$1\cdot a = a$ for all~$a$.

\begin{observation}[SAT-verified]\label{obs:ms}
$\MS(1) = 3$ and $\MS(2) = 31$.
Equivalently, $\{2,3\}$ can be 1-colored avoiding multiplicative
triples, but $\{2,3,4\}$ cannot (since $2\cdot 2=4$).
For $k=2$, $\{2,\ldots,31\}$ can be 2-colored, but
$\{2,\ldots,32\}$ cannot.
\end{observation}

\noindent
Additionally, SAT verification for $k=3$ yields
$\MS(3)=16383=2^{14}-1$ (with $\{2,\ldots,16383\}$ admitting a valid
3-coloring but $\{2,\ldots,16384\}$ not).
This suggests the formula
\begin{equation}\label{eq:ms-formula}
\MS(k) = 2^{(3^k+1)/2} - 1,
\end{equation}
which gives $\MS(1)=2^2-1=3$, $\MS(2)=2^5-1=31$, $\MS(3)=2^{14}-1=16383$.
We emphasize that \eqref{eq:ms-formula} is a \emph{conjecture}
extrapolated from three data points, not a theorem.
The exponents $2,5,14$ satisfy the recurrence $f(k)=3f(k-1)-1$ with
$f(1)=2$.

Multiplicative Schur triples have been studied recently by
Mattos, Mergoni Cecchelli, and Parczyk~\cite{mattos2024}, who
investigate the equation $xy=z$ in the integers and obtain asymptotic
density bounds.
Our exact small-$k$ values may complement their asymptotic results,
but we have not verified whether the values $\MS(1)$ and $\MS(2)$
appear in prior work.

\subsection{Mixed additive-multiplicative constraints}

One can also forbid \emph{both} additive Schur triples
$\{a,b,a{+}b\}$ and multiplicative triples $\{a,b,ab\}$
simultaneously.
We find that the resulting mixed number equals $S(k)+1$ for
$k=1,2,3$: the additive constraint completely dominates.
This is consistent with the much larger values of $\MS(k)$ compared
to $S(k)+1$ (e.g., $\MS(2)=31$ vs.\ $S(2)+1=5$).

%% ================================================================
\section*{Acknowledgments}
%% ================================================================

Computational analysis performed with AI assistance (Claude,
Anthropic).
All SAT-based results verified independently using CaDiCaL.
The Lean~4 formalization of $\DS(2,\tfrac{1}{2})=5$ was
machine-checked with zero sorry statements.
Source code and Lean proofs are available at~\cite{code-repo}.

%% ================================================================
\begin{thebibliography}{9}
%% ================================================================

\bibitem{green-ruzsa}
B.~Green and I.~Ruzsa,
``Sum-free sets in abelian groups,''
\emph{Israel J.\ Math.} \textbf{147} (2005), 157--188.

\bibitem{mattos2024}
L.~Mattos, T.~Mergoni Cecchelli, and J.~Parczyk,
``On product Schur triples in the integers,''
\emph{SIAM J.\ Discrete Math.} (2024/2025).
\url{https://epubs.siam.org/doi/10.1137/24M1632875}

\bibitem{schur-survey}
W.~Deuber, ``Developments based on Rado's dissertation
\emph{Studien zur Kombinatorik},''
in \emph{Surveys in Combinatorics} (1995), 52--74.

\bibitem{code-repo}
A.~Towell,
``Erd\H{o}s problems: computational experiments,''
2026.
Zenodo code deposit.
DOI: \texttt{10.5281/zenodo.18638635}.

\end{thebibliography}

\end{document}
